\documentclass[10pt]{article}
\usepackage[T1]{fontenc}
\usepackage[utf8]{inputenc}
\usepackage[french]{babel}
\usepackage{amsmath,amssymb,mathrsfs,mathpazo}
\usepackage{geometry}
\geometry{a4paper,margin=2cm}
\usepackage{graphicx,array,tabularx,tabularray,colortbl,multirow,multicol,booktabs,enumitem,pifont,ulem,fancybox,diagbox,cancel}
\usepackage{pstricks,pstricks-add,pst-plot,pst-tree,pst-node,pst-text,pst-func,pst-math,pst-bspline,pst-3dplot,pst-eucl,pst-solides3d}
\let\pscalculate\undefined
\usepackage{pst-bezier}
\usepackage{tikz,pgfplots,tkz-tab}
\usetikzlibrary{arrows,arrows.meta,patterns,calc,positioning,trees,decorations.markings,shapes.geometric,shapes.misc,shapes.arrows,intersections,decorations.pathmorphing,decorations.pathreplacing,backgrounds,shadows,babel}
\pgfplotsset{compat=1.18}
\usepgfplotslibrary{fillbetween}
\usepackage[np]{numprint}
\usepackage{listings,makecell,pas-tableur,twoopt}
\let\barre\undefined
\usepackage[pstricks]{tablvar}
\definecolor{darkgreen}{RGB}{0,100,0}
\providecommand{\red}{\color{red}}
\providecommand{\blue}{\color{blue}}
\providecommand{\green}{\color{green}}
\providecommand{\black}{\color{black}}
\providecommand{\white}{\color{white}}
\providecommand{\gray}{\color{gray}}
\providecommand{\cyan}{\color{cyan}}
\providecommand{\magenta}{\color{magenta}}
\providecommand{\yellow}{\color{yellow}}
\providecommand{\violet}{\color{violet}}
\providecommand{\orange}{\color{orange}}
\providecommand{\euro}{\texteuro}
\providecommand{\EUR}{\texteuro}
\providecommand{\ieme}{\textsuperscript{e}}
\providecommand{\degres}{\ensuremath{^\circ}}
\providecommand{\degree}{\ensuremath{^\circ}}
\providecommand{\vv}[1]{\overrightarrow{#1}}
\providecommand{\Coord}[1]{#1}
\providecommand{\decosix}{}
\providecommand{\starredbullet}{\ensuremath{\bullet}}
\pagestyle{empty}
\setlength{\parindent}{0pt}
\usepackage{adjustbox}\DeclareUnicodeCharacter{2236}{:}\DeclareUnicodeCharacter{2190}{\ensuremath{\leftarrow}}
\providecommand{\R}{\mathbb{R}}
\providecommand{\N}{\mathbb{N}}
\providecommand{\Z}{\mathbb{Z}}
\providecommand{\Q}{\mathbb{Q}}
\providecommand{\D}{\mathbb{D}}
\providecommand{\C}{\mathbb{C}}
\providecommand{\e}{{\mathrm{e}}}
\providecommand{\vect}[1]{\overrightarrow{#1}}
\providecommand{\vectt}[1]{\overrightarrow{\mathrm{#1}}}
\providecommand{\vvt}[1]{\overrightarrow{\mathrm{#1}}}
\providecommand{\np}[2][]{{#2~#1}}
\providecommand{\vv}[1]{{\overrightarrow{#1}}}
\renewcommand{\d}{{\mathrm{d}}}
\providecommand{\degres}{{^\circ}}
\providecommand{\degre}{{^\circ}}
\providecommand{\blue}{\color{blue}}
\providecommand{\red}{\color{red}}
\providecommand{\enskip}{\quad}
\providecommand{\llbracket}{\lbrack\!\lbrack}
\providecommand{\rrbracket}{\rbrack\!\rbrack}
\providecommand{\euro}{{\text{€}}}
\providecommand{\no}{{\text{n°}}}
\providecommand{\interval}[2]{{[#1;#2]}}
\providecommand{\intervFO}[2]{{[#1;#2[}}
\providecommand{\dotfill}{\ldots}

\providecommand{\ds}{\displaystyle}
\providecommand{\textding}[1]{#1}
\providecommand{\Oij}{\(\mathrm{O};\vec{\imath},\vec{\jmath}\)}
\providecommand{\Oijk}{\(\mathrm{O};\vec{\imath},\vec{\jmath},\vec{k}\)}
\providecommand{\Ouv}{\(\mathrm{O};\vec{u},\vec{v}\)}
\providecommand{\intervFF}[2]{\ensuremath{[#1;#2]}}
\renewcommand{\euro}{\ifmmode\text{\texteuro}\else\texteuro\fi}
\DeclareUnicodeCharacter{00B0}{\ifmmode^{\circ}\else\textdegree\fi}
\DeclareUnicodeCharacter{20AC}{\ifmmode\text{\texteuro}\else\texteuro\fi}
% Énoncés : Éducation nationale. Transcription et corrigés : APMEP, Denis Vergès et contributeurs. Découpage et classement : CoopMaths, Sébastien Lozano, Éric Elter et contributeurs. Adaptation de la présentation : Maths Explorer.
% CC BY-SA : https://coopmaths.fr/www/about/
% https://coopmaths.fr/alea/static/bac/2022/tex/bac_2022_09_sujet1_metropole_1_cor.tex
\begin{document}
\noindent\textit{APMEP / CoopMaths. CC BY-SA. Présentation adaptée par Maths Explorer.}\par\medskip
\medskip





\medskip

\begin{enumerate}
\item On considère la fonction $g$ définie sur $\R$ par : 

\[ g(x) = \dfrac{2\text{e}^x}{\text{e}^x +1}.\]

\smallskip









On a pour tout réel $x$, \: $g(x) = \dfrac{\text{e}^x \times 2 }{\text{e}^x\left(1 + \text{e}^{-x}\right)} = \dfrac{2}{1 + \text{e}^{-x}}$.

Or $\displaystyle\lim_{x \to + \infty} \text{e}^{-x} = 0$, donc $\displaystyle\lim_{x \to + \infty} 1 + \text{e}^{-x} = 1$ et enfin $\displaystyle\lim_{x \to + \infty} g(x) = \dfrac{2}{1} = 2$. \hfill Réponse \textbf{b.}

\item La dérivée seconde $f''$ est positive sur les intervalles $]- \infty~;~-1[$ et $[2~;~+ \infty[$, donc la fonction $f$ est convexe sur ces intervalles. \hfill Réponse \textbf{c.}

\item 









On a pour tout naturel $n$ :

$b_{n+1} = a_{n+1} - 2 = 0,5a_n + 1 - 2 = 0,5a_n - 1 = 0,5\left(a_n - 2 \right) = 0,5b_n$.

L'égalité $b_{n+1} = 0,5b_n$ montre que la suite $\left(b_n\right)$ est une suite géométrique de raison 0,5. 

\hfill Réponse \textbf{d.}

\item 











$\bullet~~$Comme $- 1 < \dfrac{1}{4} < 1$, on sait que $\displaystyle\lim_{n \to + \infty}\left(\dfrac14\right)^n = 0$ et donc $\displaystyle\lim_{n \to + \infty}1 + \left(\dfrac14\right)^n = 1$ ;

$\bullet~~$On a $\dfrac{n}{n + 1} = \dfrac{1}{1 + \frac1 n}$.
Or $\displaystyle\lim_{n \to + \infty} \frac1 n = 0 \Rightarrow \displaystyle\lim_{n \to + \infty} \dfrac{n}{n + 1} = 1$ et $\displaystyle\lim_{n \to + \infty}2 - \dfrac{n}{n + 1} = 1$.



On a donc  $1  + \left(\dfrac14\right)^n \leqslant u_n \leqslant 2 - \dfrac{n}{n+1}$ avec  $\displaystyle\lim_{n \to + \infty}1 + \left(\dfrac14\right)^n = \displaystyle\lim_{n \to + \infty}2 - \dfrac{n}{n + 1} = 1$. 

D'après le théorème des gendarmes, on peut dire que la suite $(u_n)$ est convergente et que $\displaystyle\lim_{n \to + \infty} u_n=1$.\hfill
Réponse \textbf{b.}

\item Soit $f$ la fonction définie sur $]0~;~+\infty[$ par $f(x) = x^2 \ln x$.










Si $F(x) = \dfrac13 x^3 \left(\ln x - \dfrac13 \right)$, alors 
$F'(x) = x^2\left(\ln x - \dfrac13 \right) + \dfrac13 x^3 \times \dfrac{1}{x} = x^2 \ln x - \dfrac13 x^2  + \dfrac13 x^2 = x^2 \ln x = f(x)$. \hfill Réponse \textbf{a.}

\item 








$2 + \dfrac{3\text{e}^{-x} - 5}{\text{e}^{-x} + 1} = \dfrac{2\text{e}^{-x} + 2  + 3\text{e}^{-x} - 5}{\text{e}^{-x} + 1} = \dfrac{5\text{e}^{-x} - 3}{\text{e}^{-x} + 1}$ ; en multipliant chaque terme par $\text{e}^x$, on obtient : $\dfrac{5 - 3\text{e}^x}{1 + \text{e}^x}$. \hfill Réponse \textbf{a.}
\end{enumerate}

\bigskip

\end{document}
